
\PassOptionsToPackage{unicode}{hyperref}
\PassOptionsToPackage{hyphens}{url}
\documentclass[
]{article}
\usepackage{xcolor}
\usepackage{amsmath,amssymb}
\setcounter{secnumdepth}{-\maxdimen} % remove section numbering
\usepackage{iftex}
\ifPDFTeX
  \usepackage[T1]{fontenc}
  \usepackage[utf8]{inputenc}
  \usepackage{textcomp} % provide euro and other symbols
\else % if luatex or xetex
  \usepackage{unicode-math} % this also loads fontspec
  \defaultfontfeatures{Scale=MatchLowercase}
  \defaultfontfeatures[\rmfamily]{Ligatures=TeX,Scale=1}
\fi
\usepackage{lmodern}
\ifPDFTeX\else
  % xetex/luatex font selection
\fi
% Use upquote if available, for straight quotes in verbatim environments
\IfFileExists{upquote.sty}{\usepackage{upquote}}{}
\IfFileExists{microtype.sty}{% use microtype if available
  \usepackage[]{microtype}
  \UseMicrotypeSet[protrusion]{basicmath} % disable protrusion for tt fonts
}{}
\makeatletter
\@ifundefined{KOMAClassName}{% if non-KOMA class
  \IfFileExists{parskip.sty}{%
    \usepackage{parskip}
  }{% else
    \setlength{\parindent}{0pt}
    \setlength{\parskip}{6pt plus 2pt minus 1pt}}
}{% if KOMA class
  \KOMAoptions{parskip=half}}
\makeatother
\usepackage{graphicx}
\makeatletter
\newsavebox\pandoc@box
\newcommand*\pandocbounded[1]{% scales image to fit in text height/width
  \sbox\pandoc@box{#1}%
  \Gscale@div\@tempa{\textheight}{\dimexpr\ht\pandoc@box+\dp\pandoc@box\relax}%
  \Gscale@div\@tempb{\linewidth}{\wd\pandoc@box}%
  \ifdim\@tempb\p@<\@tempa\p@\let\@tempa\@tempb\fi% select the smaller of both
  \ifdim\@tempa\p@<\p@\scalebox{\@tempa}{\usebox\pandoc@box}%
  \else\usebox{\pandoc@box}%
  \fi%
}
% Set default figure placement to htbp
\def\fps@figure{htbp}
\makeatother
\setlength{\emergencystretch}{3em} % prevent overfull lines
\providecommand{\tightlist}{%
  \setlength{\itemsep}{0pt}\setlength{\parskip}{0pt}}
\usepackage{bookmark}
\IfFileExists{xurl.sty}{\usepackage{xurl}}{} % add URL line breaks if available
\urlstyle{same}
\hypersetup{
  hidelinks,
  pdfcreator={LaTeX via pandoc}}

\author{}
\date{}

\begin{document}

\textbf{HYPERTRACK}

\textbf{Hyperhidrosis Monitoring and Management System}

\textbf{PROJECT PROPOSAL \textbar{} COURSE UCS503}

\textbf{Submitted to: Dr. Jeelani Asif}

\textbf{Submitted by}

\textbf{Aditya Rajput (1024240145) · AIML}

\textbf{Prabhkirat Kaur (1024240024) · AIML}

\emph{August 2026}

\textbf{Executive Overview}

HyperTrack is a proposed web-based monitoring and management system for
recording, organizing, visualizing, and analyzing user-reported sweating
episodes associated with hyperhidrosis. It is designed primarily as a
Software Engineering project, with an optional AIML component for
personalized pattern analysis and severity prediction. The system does
not diagnose or treat hyperhidrosis; it helps users organize their own
data and generate reports that can support conversations with healthcare
professionals.

\includegraphics[width=6.27083in,height=3.49765in,alt={8c60a8794a867875a246412ff0e3a2af.png}]{./media/image1.png}

\emph{Figure 1. Problem context and motivation for structured
monitoring.}

\textbf{1. Problem Statement}

Hyperhidrosis involves excessive sweating beyond normal temperature
regulation, often affecting the palms, feet, underarms, and face. A key
challenge is the lack of consistent records covering severity, duration,
affected area, triggers, activity, and stress. HyperTrack addresses this
gap by converting individual episodes into structured records that can
be analyzed over time to reveal patterns that would otherwise be missed.

\textbf{2. Objectives}

The project aims to build a reliable, user-friendly system for recording
and analyzing sweating episodes. Specific objectives include:

\begin{itemize}
\item
  Provide a user-friendly interface for recording episodes and
  associated context.
\item
  Store symptom records in a structured, secure database.
\item
  Analyze frequency, severity, duration, affected body areas, and
  possible triggers.
\item
  Deliver dashboards with graphs, statistics, and historical trends.
\item
  Support treatment/session tracking without issuing medical
  prescriptions.
\item
  Generate an organized report summarizing recorded history.
\item
  Investigate machine learning for personalized severity prediction.
\item
  Apply sound Software Engineering practice across the development
  lifecycle.
\end{itemize}

\textbf{3. Proposed Solution}

HyperTrack is proposed as a web application comprising user management,
episode recording, structured data storage, an analytics module, an
interactive dashboard, treatment/session tracking, report generation,
and an optional machine learning module. The system is positioned as a
monitoring and decision-support tool rather than a diagnostic system.

\includegraphics[width=6.27083in,height=2.60084in,alt={05156ca988e876ab6126ac51f2b3e0eb.png}]{./media/image2.png}

\emph{Figure 2. Eight functional modules forming the core application.}

\textbf{4. System Architecture}

The application follows a three-tier architecture. The presentation
layer (HTML, CSS, JavaScript) delivers the user interface; the
application layer (Python and Flask) exposes REST APIs and hosts
analytics and machine learning services; and the data layer (MySQL)
stores users, episodes, triggers, and treatment records. This separation
keeps the system easier to test, maintain, and scale.

\includegraphics[width=6.27083in,height=2.6577in,alt={5365f1632a8974b5bc51e010177f7b7a.png}]{./media/image3.png}

\emph{Figure 3. Three-tier architecture with analytics and ML services.}

\textbf{5. System Workflow}

The core workflow begins with user login, followed by episode recording,
input validation, and database storage. Historical records then pass
through the analytics module to populate the dashboard, after which the
optional machine learning module can add a severity insight before the
user generates a structured monitoring report.

\includegraphics[width=6.27083in,height=2.4872in,alt={6ef08654f60bc9433bc0bc45c9db534a.png}]{./media/image4.png}

\emph{Figure 4. End-to-end operational workflow.}

\textbf{6. Episode Data Collection}

Each recorded episode captures the date and time, sweating severity,
duration, affected body area, activity, self-reported stress level,
possible trigger, temperature and humidity where available, and any
additional notes. This structured record forms the basis for both
statistical analysis and future machine learning functionality.

\includegraphics[width=6.27083in,height=3.25221in,alt={7301f7432e847e18dc0c5fead7fde1be.png}]{./media/image5.png}

\emph{Figure 5. Structured data captured for each episode.}

\textbf{7. Data Processing and Analytics}

Stored records are cleaned and processed to calculate daily and weekly
frequency, average severity, average duration, frequently affected
areas, commonly recorded triggers, severity trends, and associations
between stress and severity. Results are presented as patterns or
associations in the user\textquotesingle s recorded data rather than
medical conclusions.

\includegraphics[width=6.27083in,height=1.65855in,alt={3c12c8a8d5b2e6c1de53d6db4292637f.png}]{./media/image6.png}

\emph{Figure 6. Analytics pipeline and derived metrics.}

\textbf{8. Machine Learning Component}

As the AIML component, the project investigates supervised severity
classification using features such as stress, temperature, humidity,
activity, time of day, previous severity, duration, caffeine intake, and
affected body area. Candidate algorithms include Logistic Regression,
Decision Trees, and Random Forest, with Random Forest considered the
initial model for its ability to handle multiple features and report
feature importance. The output is a severity category of Low, Moderate,
or High, offered strictly as a supplementary insight.

\includegraphics[width=6.27083in,height=1.67029in,alt={dd7ad37146d69adea1d5b481f2485a57.png}]{./media/image7.png}

\emph{Figure 7. Proposed machine learning pipeline.}

\textbf{9. Testing and Validation}

Testing runs throughout development. Unit testing covers individual
functions such as login, episode creation, and validation; integration
testing verifies communication across the frontend, backend, database,
and analytics modules; system testing validates the complete workflow
from login through report generation; and user acceptance testing
evaluates usability and clarity with selected users. If implemented, the
ML module is evaluated using accuracy, precision, recall, F1-score, and
a confusion matrix.

\includegraphics[width=6.27083in,height=1.67029in,alt={e7b2ddc87890a40a37b374dfaf7459ad.png}]{./media/image8.png}

\emph{Figure 8. Multi-level testing and validation strategy.}

\textbf{10. Technology Stack}

The proposed stack relies primarily on free and open-source
technologies: HTML/CSS/JavaScript for the frontend; Python, Flask, and
REST APIs for the backend; MySQL, Pandas, and NumPy for data;
Scikit-learn, Chart.js, and Matplotlib for machine learning and
visualization; and Git/GitHub with VS Code or PyCharm for tooling and
version control.

\includegraphics[width=6.27083in,height=3.31131in,alt={0307d3ae3bd9f5084f0054af31b9f2a4.png}]{./media/image9.png}

\emph{Figure 9. Technology stack across the application.}

\textbf{11. Development Timeline}

Development spans a 15-week semester using an iterative approach,
prioritizing core monitoring functionality before advanced analytics and
machine learning features. A buffer period is maintained during the
final development and testing stages to absorb unexpected technical
issues.

\includegraphics[width=6.27083in,height=2.08in,alt={a0db15f41a6fb1a89481732076611b2e.png}]{./media/image10.png}

\emph{Figure 10. Semester development roadmap.}

\textbf{12. Resources and Budget}

The project is developed by two students. Aditya Rajput (AIML) leads
backend development, database integration, machine learning, and system
integration. Prabhkirat Kaur (AIML) leads frontend development, UI/UX,
documentation, testing, and analytics integration. Both members
participate in requirements analysis, system design, testing, and final
documentation.

The software stack is built entirely on freely available or open-source
tools --- Python, Flask, MySQL, Pandas, NumPy, Scikit-learn,
HTML/CSS/JavaScript, Git/GitHub, and VS Code or PyCharm --- so no
external funding is required for the initial implementation. Future
cloud hosting or external services, if needed, may introduce additional
costs.

\textbf{13. Risks and Mitigation}

Key risks include limited data availability for the ML prototype,
protection of sensitive health data, the possibility of medical
misinterpretation, ML reliability, development delays, and user
adoption. Each is addressed with a concrete mitigation strategy,
summarized below.

\includegraphics[width=6.27083in,height=2.87823in,alt={68cbd3e5834a6f0cd49f40c28a4a332c.png}]{./media/image11.png}

\emph{Figure 11. Key project risks and mitigation strategies.}

\textbf{14. Expected Outcomes}

At completion, HyperTrack is expected to be a functional web application
that lets users maintain structured episode records, understand
frequency and severity trends, identify common triggers, track treatment
history, and generate an organized monitoring report --- with optional
ML-based severity insight. From a Software Engineering standpoint, the
project demonstrates requirements engineering, system architecture,
database management, modular development, testing, and iterative
delivery.

\includegraphics[width=6.27083in,height=2.31442in,alt={ba6138c4416f03f5d67e1921ef8d486d.png}]{./media/image12.png}

\emph{Figure 12. Expected project outcomes.}

\textbf{Clinical Scope Note}

HyperTrack is a monitoring and decision-support system. Its analytics
and ML outputs are supplementary insights drawn from user-recorded data
--- they are not medical diagnosis, treatment recommendations, or
medical advice.

Aditya Rajput --- Roll No. 1024240145 --- arajput\_be24@thapar.edu

Prabhkirat Kaur --- Roll No. 1024240024 --- pkaur2\_be24@thapar.edu

\end{document}
